% ============================================================
%  template.tex — Notas Acadêmicas
%  Autor: <Gabriel de Souza Dourado>
% ============================================================

\documentclass[12pt, a4paper]{article}

% ── Codificação e Língua ────────────────────────────────────
\usepackage[utf8]{inputenc}
\usepackage[T1]{fontenc}
\usepackage[brazil]{babel}

% ── Tipografia ──────────────────────────────────────────────
\usepackage{lmodern}
\usepackage{microtype}

% ── Margens ─────────────────────────────────────────────────
\usepackage[
  top=2.5cm, bottom=2.5cm,
  left=2.8cm, right=2.8cm
]{geometry}
\setlength{\skip\footins}{1.5em}

% ── Matemática ──────────────────────────────────────────────
\usepackage{amsmath, amssymb, amsthm}
\usepackage{mathtools}

% ── Algoritmos e Código ─────────────────────────────────────
\usepackage{algorithm}
\usepackage{algpseudocode}
\usepackage{listings}
\usepackage{xcolor}

% ── Pseudocódigo ─────────────────────────────────────
\renewcommand{\algorithmicrequire}{\textbf{Entrada:}}
\renewcommand{\algorithmicensure}{\textbf{Saída:}}

% ── Figuras e Tabelas ───────────────────────────────────────
\usepackage{graphicx}
\usepackage{booktabs}
\usepackage{caption}
\usepackage{float}

% ── Hiperlinks ──────────────────────────────────────────────
\usepackage[
  colorlinks=true,
  linkcolor=blue!70!black,
  citecolor=blue!70!black,
  urlcolor=blue!70!black
]{hyperref}

% ── Cabeçalho e Rodapé ──────────────────────────────────────
\usepackage{fancyhdr}
\pagestyle{fancy}
\fancyhf{}
\fancyhead[L]{\small\textit{\docsubtitle}}
\fancyhead[R]{\small\textit{\docauthor}}
\fancyfoot[C]{\small\thepage}
\renewcommand{\headrulewidth}{0.4pt}

% ── Cores ───────────────────────────────────────────────────
\definecolor{mainblue}{RGB}{30, 80, 160}
\definecolor{codegray}{RGB}{245, 245, 245}
\definecolor{codecomment}{RGB}{100, 100, 100}
\definecolor{codestring}{RGB}{163, 21, 21}
\definecolor{codekeyword}{RGB}{30, 80, 160}

% ── Estilo de Código (C++) ───────────────────────────────────
\lstdefinestyle{cppstyle}{
  language=C++,
  backgroundcolor=\color{codegray},
  basicstyle=\ttfamily\small,
  keywordstyle=\color{codekeyword}\bfseries,
  commentstyle=\color{codecomment}\itshape,
  stringstyle=\color{codestring},
  numberstyle=\tiny\color{codecomment},
  numbers=left,
  numbersep=8pt,
  breaklines=true,
  frame=single,
  framerule=0pt,
  rulecolor=\color{codegray},
  showstringspaces=false,
  tabsize=2,
  captionpos=b,
}
\lstset{style=cppstyle}

% ── Ambientes de Teorema ─────────────────────────────────────
\theoremstyle{definition}
\newtheorem{definition}{Definição}[section]
\newtheorem{example}{Exemplo}[section]

\theoremstyle{plain}
\newtheorem{theorem}{Teorema}[section]
\newtheorem{lemma}[theorem]{Lema}
\newtheorem{corollary}[theorem]{Corolário}
\newtheorem{proposition}[theorem]{Proposição}

\theoremstyle{remark}
\newtheorem*{remark}{Observação}
\newtheorem*{note}{Nota}

% ── Espaçamento ─────────────────────────────────────────────
\usepackage{setspace}
\setstretch{1.3}

% ============================================================
%  METADADOS — altere apenas estas 4 linhas em cada documento
% ============================================================
\newcommand{\docauthor}{Gabriel de Souza Dourado}
\newcommand{\doctitle}{Teoria da Computação}
\newcommand{\docsubtitle}{Gramáticas Regulares}
\newcommand{\docdate}{\today}
% ============================================================

\begin{document}

% ── Título ───────────────────────────────────────────────────
\begin{center}
  {\large \textcolor{mainblue}{\textsc{\docsubtitle}}}\\[0.4em]
  {\LARGE \textbf{\doctitle}}\\[0.6em]
  {\small \docauthor \quad·\quad \docdate}
  \vspace{0.4em}
  \hrule
\end{center}

\vspace{1em}

% ── Sumário (descomente se quiser) ───────────────────────────
% \tableofcontents
% \newpage

% ------------------------------------------------------------
\section{Introdução}
% ------------------------------------------------------------
Uma \textbf{Gramática de Chomsky} consiste em um conjunto de regras de produção que, quando aplicadas de forma sucessiva, geram palavras. O conjunto de todas as palavras que podem ser geradas por uma gramática define uma linguagem.

\section{Definição}

\begin{definition}[Gramática de Chomsky]
Uma \textbf{Gramática de Chomsky} é definida por uma 4-tupla ordenada:

\[
G = (V, T, P, S)
\]

onde:

\begin{itemize}
    \item \(V\) é um conjunto finito de símbolos não-terminais.
    \item \(T\) é um conjunto finito de símbolos terminais, com \(V \cap T = \emptyset\).
    \item \(P \subseteq (V \cup T)^{+} \times (V \cup T)^*\) é um conjunto finito de produções.
    \item \(S \in V\) é o símbolo inicial da gramática.
\end{itemize} 
\end{definition}
    
\section{Gramáticas Regulares}

Uma gramática é \textbf{Regular} se atender, e unicamente atender, a uma das seguintes restrições:

\begin{definition}[Gramática Linear à Direita]
Dado $A, B \in V$ e $\omega \in T^*$, seja $G = (V, T, P, S)$, se toda produção tem a forma:
\[
A \rightarrow \omega B \quad \text{ou} \quad A \rightarrow \omega
\]
ela é dita uma Gramática Linear à Direita. 
\end{definition}

\begin{definition}[Gramática Linear à Esquerda]
Dado $A, B \in V$ e $\omega \in T^*$, seja $G = (V, T, P, S)$, se toda produção tem a forma:
\[
A \rightarrow B \omega \quad \text{ou} \quad A \rightarrow \omega
\]
ela é dita uma Gramática Linear à Esquerda. 
\end{definition}


\begin{definition}[Gramática Linear Unitária à Direita]
Dado $A, B \in V$ e $\omega \in T^*$, seja $G = (V, T, P, S)$, se toda produção tem a forma:
\[
A \rightarrow \omega B \quad \text{ou} \quad A \rightarrow \omega, \quad |\omega| \le 1
\]
ela é dita uma Gramática Linear Unitária à Direita. 
\end{definition}

\begin{definition}[Gramática Linear Unitária à Esquerda]
Dado $A, B \in V$ e $\omega \in T^*$, seja $G = (V, T, P, S)$, se toda produção tem a forma:
\[
A \rightarrow B \omega \quad \text{ou} \quad A \rightarrow \omega, \quad |\omega| \le 1
\]
ela é dita uma Gramática Linear Unitária à Esquerda. 
\end{definition}

\section{Conversão de GLUD para AFND-$\varepsilon$}

Vamos converte uma \textbf{Gramática Linear Unitária a Direita (GLUD)} em um \textbf{Autômato Finito Não Determinístico com transições vazias (AFND-$\epsilon$)}.

\subsection{Descrição Geral}

Dada uma GLUD $G = (V, T, P, S)$, deseja-se construir um AFND-$\epsilon$ $A = (Q, \Sigma, \delta, q_{0}, F)$. Por definição, consideramos as seguintes igualdades:
\begin{center}
    $Q = V \cup \{q_f\}$, $\Sigma = T$ $q_{0} = S$ e $F = \{q_f\}$.
\end{center}
A função $\delta$ de transição é definida da seguinte forma: \\
\begin{center}
\begin{tabular}{c|c}
  \textbf{Regra de Produção} & \textbf{Função de Transição $\delta$} \\\hline
  $A \rightarrow \epsilon$ & $\delta(A, \epsilon) = \{q_f\}$ \\
  $A \rightarrow \sigma$ & $\delta(A, \sigma) = \{q_f\}$ \\
  $A \rightarrow B$ & $\delta(A, \epsilon) = \{B\}$ \\
  $A \rightarrow \sigma B$ & $\delta(A, \sigma) = \{B\}$ \\
\end{tabular}
\end{center}

\subsection{Algoritmo}
O procedimento do algoritmo pode ser descrito da seguinte forma:

\begin{algorithm}[H]
  \caption{$\textsc{Converte}(G)$}
  \label{alg:converte}
  \begin{algorithmic}[1]
    \Require Uma GLUD $G = (V, T, P, S)$
    \Ensure Um AFND-$\varepsilon$ $M = (Q, \Sigma, \delta, q_0, F)$,
             onde $\delta : Q \times (\Sigma \cup \{\varepsilon\}) \rightharpoonup 2^Q$

    \State $q_f \leftarrow$ novo estado; \quad $Q \leftarrow V \cup \{q_f\}$
    \State $\Sigma \leftarrow T$
    \State $q_0 \leftarrow S$
    \State $F \leftarrow \{q_f\}$
    \State $\delta \leftarrow \emptyset$

    \For{cada produção $p \in P$}
      \If{$p$ é da forma $A \rightarrow aB$}
        \State $\delta(A,\, a) \leftarrow \delta(A,\, a) \cup \{B\}$
      \ElsIf{$p$ é da forma $A \rightarrow a$}
        \State $\delta(A,\, a) \leftarrow \delta(A,\, a) \cup \{q_f\}$
      \ElsIf{$p$ é da forma $A \rightarrow B$}
        \State $\delta(A,\, \varepsilon) \leftarrow \delta(A,\, \varepsilon) \cup \{B\}$
      \ElsIf{$p$ é da forma $A \rightarrow \varepsilon$}
        \State $\delta(A,\, \varepsilon) \leftarrow \delta(A,\, \varepsilon) \cup \{q_f\}$
      \EndIf
    \EndFor

    \State \Return $M = (Q,\, \Sigma,\, \delta,\, q_0,\, F)$
  \end{algorithmic}
\end{algorithm}

% ============================================================
%  REFERÊNCIAS
% ============================================================
\newpage
\begin{thebibliography}{9}

\bibitem{menezes}
MENEZES, P. B.\@.
\textit{Linguagens Formais e Autômatos}.
6.~ed. Porto Alegre: Bookman, 2011.

\bibitem{hopcroft}
HOPCROFT, J. E.; MOTWANI, R.; ULLMAN, J. D.\@.
\textit{Introduction to Automata Theory, Languages, and Computation}.
3.~ed. Boston: Pearson, 2006.

\end{thebibliography}

\end{document}